% !TeX root = ../main.tex

\sysusetup{
  keywords  = {知识管理系统，检索增强生成，混合检索，访问控制，企业文档},
  keywords* = {knowledge management system, retrieval-augmented generation, hybrid retrieval, access control, enterprise documents},
  keywords** = {gestion des connaissances, generation augmentee par recherche, recherche hybride},
}

\begin{abstract}
  面向企业私有文档场景，传统知识管理系统普遍存在检索入口分散、文档权限约束不足以及自然语言交互能力有限等问题；仅依赖大语言模型参数记忆的问答方式，又难以满足企业知识服务对可追溯性、时效性和访问控制的要求。针对上述问题，本文设计并实现了一套基于检索增强生成的智能知识管理系统，将用户认证、组织标签管理、文档接入、知识入库、混合检索与流式问答纳入统一的软件架构。系统以后端服务、对象存储、缓存、搜索引擎与大模型接口协同构建完整处理链路：在文档接入阶段支持分片上传、断点续传、对象合并、文本解析与向量化写入；在问答阶段结合组织标签、文档公开属性与用户身份执行权限约束下的混合检索，并基于检索结果组织提示词，实现带来源引用的对话生成。

  论文围绕需求分析、总体设计、关键模块实现与系统测试展开论述，对权限模型、异步入库链路、检索流程和会话编排机制进行了细化说明。系统测试以第二章需求为主线，对身份建立、文档接入、混合检索、权限控制、文件服务和流式交互进行了联合验证。结果表明，500 份测试文档均完成上传、续传与对象合并，100 份全链路样本中有 99 份完成解析与索引写入；40 条检索查询全部在 Top5 内命中目标文档，25 组权限矩阵检查未出现越权返回，37 项文件服务检查全部通过。与此同时，系统已经具备流式回答与历史记录读取能力，但多轮历史追问仍存在语义偏移。研究结果说明，将检索增强生成能力嵌入具有权限控制与文档生命周期管理能力的软件系统，能够为企业内部知识服务提供一条可实现、可扩展的工程路径。

\end{abstract}

\begin{abstract*}
  In enterprise scenarios built around private documents, conventional knowledge management systems often suffer from fragmented retrieval entry points, weak access control over documents, and limited natural-language interaction. Question answering based solely on the parametric memory of large language models also struggles to satisfy the requirements of traceability, timeliness, and permission enforcement in enterprise knowledge services. To address these issues, this thesis designs and implements an intelligent knowledge management system based on retrieval-augmented generation. The system integrates user authentication, organizational tag management, document ingestion, knowledge indexing, hybrid retrieval, and streaming dialogue into a unified software architecture. A complete processing pipeline is constructed through the coordination of backend services, object storage, cache, search engine, and external model APIs. During document ingestion, the system supports chunked upload, resumable transfer, object composition, text parsing, and vector indexing. During question answering, it performs permission-constrained hybrid retrieval according to user identity, organizational scope, and document visibility, and then assembles prompts from retrieved evidence to generate responses with source references.

  The thesis is organized around requirement analysis, overall design, key-module implementation, and system testing, with detailed discussions of the permission model, the asynchronous ingestion pipeline, the retrieval workflow, and the conversation orchestration mechanism. System testing is organized around the requirements defined in Chapter 2 and jointly verifies identity establishment, document ingestion, hybrid retrieval, permission enforcement, file services, and streaming interaction. The results show that all 500 test documents completed upload, resumable transfer, and object composition, while 99 of the 100 full-chain samples completed parsing and indexing. All 40 retrieval queries hit the target document within Top5, the 25 permission-matrix checks produced zero unauthorized results, and all 37 file-service checks passed. Meanwhile, streaming response and history retrieval are available, but multi-turn follow-up based on historical context still shows semantic drift. These results suggest that embedding retrieval-augmented generation into a software system equipped with access control and document lifecycle management provides a practical and extensible engineering path for enterprise knowledge services.
\end{abstract*}

% 法语（第三种语言）摘要，本科不需要就全部注释掉即可，研究生请直接注释
%\begin{abstract**}
%  La longeur de l'abstract français doit faire référence à l'abstract chinois. Le titre de l'abstract français est «ABSTRACT». La première lettre de chaque mot-clé doit être en gros, et les mots-clés doivent être séparer par des virgules et espaces. Les mots-clés de chaque langue doivent être correspondants les uns aux autres.
%\end{abstract**}
